\begin{frame}{자주 묻는 질문 — $r_t$ 와 KL}
  \small
  \begin{itemize}
    \item \kb{Q. $r_t(\theta_{\mathrm{old}})=1$ 이 항상 성립하는
      이유는?} 방금 데이터를 모은 시점에는 ``지금 정책''과
      ``데이터를 모은 정책''이 완전히 같으므로 같은 행동에 부여하는
      확률도 같다(비율이 1). 그 이후 경사 상승 스텝을 거치며
      $\theta$ 가 $\theta_{\mathrm{old}}$ 에서 멀어지면 $r_t$ 가
      1에서 벗어난다.
    \item \kb{Q. 중요도 샘플링 하나만으론 왜 부족한가?} 재가중은
      해 주되, 그 재가중이 만드는 그래디언트의 \textbf{크기에는
      아무 제약이 없다} — $r_t$ 가 매우 크면 그 한 샘플이 전체
      업데이트를 지배. 클리핑(11.2)이 바로 이 ``크기 제한''.
    \item \kb{Q. $A_t=0$ 인 샘플은 $r_t$ 가 아무리 커도 무관?}
      \textbf{맞다} — $L^{\mathrm{IS}}=r_tA_t$ 이므로 $A_t=0$ 이면
      기여는 항상 0. 단, $A_t$ 가 0에 \textit{가까운} 샘플은
      $r_t$ 가 크면 $r_tA_t$ 가 작지 않을 수 있어,
      ``약하게 좋았다/나빴다''는 신호가 극단 샘플의 노이즈에
      묻힌다.
  \end{itemize}
\end{frame}
